\SourceLeaf{397}%
\clearpage
\renewcommand{\BookRunningTitle}{On the Theory of Complex Numbers}
\pdfbookmark[0]{On the Theory of Complex Numbers}{nachlass-complexe-zahlen}
\begin{center}\Large ON THE THEORY OF COMPLEX NUMBERS.\end{center}
\BookRule{.15\textwidth}
\pdfbookmark[1]{I. New Theory of the Decomposition of Cubes}{nachlass-complexe-I}
\begin{center}[I.]\par\medskip\large NEW THEORY OF THE DECOMPOSITION OF CUBES.\end{center}

I. We assume that the equation $x^3+y^3+z^3=0$ has a solution, namely $x=a,\ y=b,\ z=c$, where $a,b,c$ have no common divisor and are therefore also pairwise coprime. We set
\[
\begin{aligned}
b+c&=\alpha\\
c+a&=\beta\\
a+b&=\gamma
\end{aligned}
\]
where $\alpha,\beta,\gamma$ must also be pairwise coprime. For if $\alpha$ and $\beta$ had a common divisor, it would also divide $a^3$ and $b^3$, so $a$ and $b$ would have to have a common divisor too.

We now have
\[
(\beta+\gamma-\alpha)^3+(\gamma+\alpha-\beta)^3+(\alpha+\beta-\gamma)^3=0
\]
but the identity
\[
\begin{aligned}
&(\beta+\gamma-\alpha)^3+(\gamma+\alpha-\beta)^3+(\alpha+\beta-\gamma)^3\\
&=(\alpha+\beta+\gamma)^3-24\alpha\beta\gamma
\end{aligned}
\]
holds. Consequently,
\[
(\alpha+\beta+\gamma)^3=24\alpha\beta\gamma
\]

\SourceLeaf{398}%
If $\alpha,\beta,\gamma$ are real numbers, then $\alpha+\beta+\gamma$ will be divisible by $3$, hence $\alpha+\beta+\gamma)^3$ by $27$, and therefore $\alpha\beta\gamma$ by $9$. Thus one of the numbers $\alpha,\beta,\gamma$, for example $\gamma$, must be divisible by $9$, and so must $c^3$; consequently $c$ is divisible by $3$.

If, however, $\alpha,\beta,\gamma$ are imaginary numbers, we conclude that $\alpha+\beta+\gamma$ must be divisible by $1-\varepsilon$, hence $24\alpha\beta\gamma$ by $(1-\varepsilon)^3$, and therefore $\alpha\beta\gamma$ by $1-\varepsilon$. Thus one of the numbers $\alpha,\beta,\gamma$ is divisible by $1-\varepsilon$, and consequently so is one of the numbers $a,b,c$.

II. In general, we have the identity
\[
\begin{aligned}
&(p+q+r)^3+(p+q\varepsilon+r\varepsilon\varepsilon)^8+(p+q\varepsilon\varepsilon+r\varepsilon)^3\\
&=27pqr+3(p+q+r)(p+q\varepsilon+r\varepsilon\varepsilon)(p+q\varepsilon\varepsilon+r\varepsilon)
\end{aligned}
\]
Consequently, if $p+q+r=0$, then
\[
(p+q\varepsilon+r\varepsilon\varepsilon)^3+(p+q\varepsilon\varepsilon+r\varepsilon)^3-27pqr=0
\]

If $p,q,r,$ are themselves cubes, namely $=a^3,b^3,c^3$ respectively, that is, if the equation $x^3+y^3+z^3=0$ has a solution, then, on setting
\[
\begin{aligned}
a^3+b^3\varepsilon+c^3\varepsilon\varepsilon&=a'\\
a^3+b^3\varepsilon\varepsilon+c^3\varepsilon&=b'\\
-3abc&=c'
\end{aligned}
\]
we also have $a'^3+b'^3+c'^3=0$. From this new solution a third can be derived in the same way, and so on. It is easy to see that if the first solution consists of real numbers, the third will too.

It should also be noted that if $a,b,c$ have no common factor, the same will hold for $a',b',c'$, apart from the factor $1-\varepsilon$. For
\[
\begin{aligned}
\frac{a'}{1-\varepsilon}&=-\varepsilon\varepsilon a^3+\varepsilon b^3=a^2-\varepsilon c^3=-b^3+\varepsilon\varepsilon c^3\\
\frac{b'}{1-\varepsilon}&=a^3-\varepsilon b^3=-\varepsilon\varepsilon a^3+\varepsilon c^3=\varepsilon\varepsilon b^3-c^3\\
\frac{c'}{1-\varepsilon}&=(\varepsilon\varepsilon-1)abc
\end{aligned}
\]
The first two numbers therefore have no factor in common with $a$, $b$, or $c$, and cannot be divisible by $1-\varepsilon$ either, unless $a,b,c$
\SourceLeaf{399}%
are all divisible by $1-\varepsilon$; thus they also have no factor in common with the third.

III. But the reverse path is also open. We have seen that one of the quantities is divisible by $1-\varepsilon$; let it be $c$. Since we can substitute $a\varepsilon$ or $a\varepsilon\varepsilon$ for $a$, and likewise $b\varepsilon$ or $b\varepsilon\varepsilon$ for $b$, we may assume that $a$ is either $\equiv1$ or $\equiv-1$. We shall assume the former, since in the other case $b\equiv1$ would hold, and it would only need to be interchanged with $a$. Accordingly, we set
\[
\begin{aligned}
a&=1+3\alpha\\
b&=-1+3\beta
\end{aligned}
\]
and
\[
\begin{aligned}
\frac{a\varepsilon+b\varepsilon\varepsilon}{\varepsilon-\varepsilon\varepsilon}&=1+(\varepsilon\varepsilon-\varepsilon)(\alpha\varepsilon+\beta\varepsilon\varepsilon)=A\\
\frac{a\varepsilon\varepsilon+b\varepsilon}{\varepsilon-\varepsilon\varepsilon}&=-1+(\varepsilon\varepsilon-\varepsilon)(\alpha\varepsilon\varepsilon+\beta\varepsilon)=B\\
\frac{a+b}{\varepsilon-\varepsilon\varepsilon}&=(\varepsilon\varepsilon-\varepsilon)(\alpha+\beta)=C
\end{aligned}
\]
where $A+B+C=0$, and $ABC=\frac{a^3+b^3}{(\varepsilon-\varepsilon\varepsilon)^3}=\left(\frac c{\varepsilon\varepsilon-\varepsilon}\right)^3$

Since here
\[
\begin{aligned}
a&=-\varepsilon A+\varepsilon\varepsilon B\\
b&=\varepsilon\varepsilon A-\varepsilon B
\end{aligned}
\]
$A$ and $B$ cannot have a common factor, because it would otherwise also be a common factor of $a$ and $b$. Since $A+B+C=0$, $C$ consequently cannot have a factor in common with either $A$ or $B$. It follows easily that $A$ and $B$, and hence also $C$, are cubes. For $\left(\frac c{\varepsilon\varepsilon-\varepsilon}\right)^3$ will be divisible by $\varepsilon-\varepsilon\varepsilon$, hence also by $(\varepsilon-\varepsilon\varepsilon)^3$, or $\alpha+\beta$ by $3$; therefore $A\equiv1$, $B\equiv-1\ (\text{mod. }3)$.

If we now set
\[
\begin{aligned}
A&=a'^3\\
B&=b'^3\\
C&=c'^3
\end{aligned}
\]
we have derived from the solution of the equation $x^3+y^3+z^3=0$

\SourceLeaf{400}%
\[
\begin{aligned}
x&=a\\
y&=b\\
z&=c
\end{aligned}
\]
another solution,
\[
\begin{aligned}
x&=a'\\
y&=b'\\
z&=c'
\end{aligned}
\]
\[
\text{where }a'^3b'^3c'^3=\frac{c^3}{(\varepsilon\varepsilon-\varepsilon)^3}
\]
so $c'$ will contain the factor $1-\varepsilon$ one fewer time than $c$. But this is absurd if $c$ is divisible only by a fixed power of $1-\varepsilon$, that is, if $c$ differs from $0$. For by continuing these operations we would otherwise eventually reach a solution in which $z$ was not divisible by $1-\varepsilon$ at all, contrary to (I).
\BookRule{.3\textwidth}

We can take a similar path for powers with exponent $5$. Namely, if $a^5+b^5+c^5=0$, we set $b+c=\alpha,\ c+a=\beta,\ a+b=\gamma$, and then
\[
\begin{aligned}
0&=(2a)^5+(2b)^5+(2c)^5\\
&=(\beta+\gamma-\alpha)^5+(\gamma+\alpha-\beta)^5+(\alpha+\beta-\gamma)^5\\
&=(\alpha+\beta+\gamma)^5-80\alpha\beta\gamma(\alpha\alpha+\beta\beta+\gamma\gamma)
\end{aligned}
\]

\Needspace{7\baselineskip}
But $(\alpha+\beta+\gamma)^5=80\alpha\beta\gamma(\alpha\alpha+\beta\beta+\gamma\gamma)$ cannot hold unless one of the numbers $\alpha,\beta,\gamma$ is divisible by $1-\varepsilon$. For if none of them were divisible, both $\alpha+\beta+\gamma$ and $\alpha\alpha+\beta\beta+\gamma\gamma$ would have to be divisible by $1-\varepsilon$, and consequently so would $2(\alpha\alpha+\beta\beta+\gamma\gamma)+2(\alpha+\beta+\gamma)(\alpha+\beta-\gamma)=(2\alpha+\beta)^2+3\beta\beta$, which is impossible.

We can also present this as follows. If $a^5+b^5+c^5=0$, then
\[
\begin{aligned}
4(a+b+c)^5&=5(b+c)(c+a)(a+b)\\
&\qquad[(a+2b+3c)^2+3(a+c)^2-8(a+b+c)c]\\
&=5(b+c)(c+a)(a+b)\\
&\qquad[(b-c)^2+3(b+c)^2+4(a+b+c)a]
\end{aligned}
\]
\[
4(a+b+c)^5+5abc[(b-c)^2+3(b+c)^2]=5(a+b+c)\{\dotsfour\}
\]
The rest of the proof would be very similar to the preceding one.

But if we try the same procedure with seventh powers, we cannot prove that for a given solution
\SourceLeaf{401}%
\[
a^7+b^7+c^7=0
\]
one of the quantities $a,b,c$ must be divisible by $7$. For we obtain only
\[
\begin{aligned}
(\alpha+\beta+\gamma)^7={}&5\alpha\beta\gamma\{3(\alpha^4+\beta^4+\gamma^4)\\
&\qquad+10(\alpha\alpha\beta\beta+\alpha\alpha\gamma\gamma+\beta\beta\gamma\gamma)\}
\end{aligned}
\]
which can hold without $\alpha,\beta,\gamma$ being divisible by $1-\varepsilon$.

Hopefully, however, this can in future be derived from the nature of the determinants and the units.
\BookRule{.3\textwidth}
